\chapter{What project finance is, and when to use it}\label{ch:2}

On April 21, 2016, SunEdison, Inc.\ filed for Chapter~11 in the US Bankruptcy Court for the Southern District of New York. Seven months earlier its consolidated balance sheet had shown total liabilities of USD~16.1~billion, of which USD~11.7~billion was debt (SunEdison 2015). It filed with commitments for up to USD~300~million of new financing to keep the business running in bankruptcy (SunEdison 2016).

Two companies that SunEdison controlled were missing from the list of debtors. TerraForm Power and TerraForm Global were listed on NASDAQ, owned operating wind and solar plants that SunEdison had built or bought, and paid their shareholders dividends from those plants' cash. SunEdison held most of their votes. Much of the USD~11.7~billion of debt on SunEdison's consolidated balance sheet was not SunEdison's to repay: it had been borrowed by the TerraForm companies and by the individual projects beneath them, and the lenders who provided it could look only to those companies and those projects. Neither TerraForm company was pulled into the bankruptcy, and in 2017 Brookfield bought control of TerraForm Power and all of TerraForm Global as going concerns (TerraForm Power 2018).

The separation did not hold cleanly. TerraForm Power's audited 2015 accounts carried a going-concern explanatory note, partly because a bankruptcy court might consolidate it with SunEdison. Its 2015 annual report was not filed until December 5, 2016, which put its bond indentures and its stock exchange listing at risk. Five months before the filing, SunEdison had used its voting control to remove TerraForm Power's chief executive and chief financial officer (TerraForm Power 2016). The plants kept generating electricity. The companies that owned them spent two years fighting their own sponsor.

Chapter~1's Llano Pardo deal used the same structure without naming it. What kept the TerraForm plants and their lenders out of SunEdison's bankruptcy, and why the companies' governance was dragged in anyway, are the first things a project financier needs to understand about it.

\section*{What you will be able to do}
\begin{itemize}
  \item Explain what project finance is, how ring-fencing and limited recourse work, and how it differs from corporate finance, leasing, reserve-based lending, acquisition finance, securitization, and structured finance (Capability 1).
  \item Compare the cost, time, and control burdens of project finance with balance-sheet borrowing for a given project, and say when project finance is the wrong tool (Capabilities 1 and 15).
  \item Place any risk on a project with the party best able to manage it, and recognize when that principle fails in practice (Capability 2).
\end{itemize}

\section{The project company and its ring-fence}\label{sec:2.1}

\subsection{A company with one asset and no history}\label{ssec:2.1.1}

When Sterrenberg Bank's credit committee approved its half of the Llano Pardo loan in 2016 (\cref{ch:1}), the borrower had no audited accounts, no revenue, no employees of note, and no track record. Llano Pardo Solar SA had been incorporated in Corredana for one purpose: to build and own a 50~MW solar plant on 152 hectares of leased plateau and sell its output to the state utility, ELNACOR, for 20 years. A corporate lender would have asked how the borrower had performed over the past five years. The committee could only ask what the plant would do over the next twenty, and who had promised what.

That shift in the question defines the discipline. \term{Project finance} is financing in which lenders are repaid from the cash flows of a single ring-fenced project and secured on its assets and contracts, with limited or no recourse to the sponsors. Each of its three features changes how a lender works. Because repayment comes from one project's cash flows, the lender studies the project, not the owner: its output, its revenue contract, its costs. Because the security is the project's assets and contracts, the lender's protection on default is the ability to take over a working plant and its contracts, not the ability to sell a pile of steel. Because recourse to the owners is limited or absent, the lender cannot fall back on a parent company's balance sheet when the forecasts prove wrong.

The legal device that makes all three features possible is the \term{special purpose vehicle} (SPV): a company created to own and operate one project and to do nothing else, so that the project's assets, contracts, and debts are held apart from its owners'. The legal form follows the host country's company law. Llano Pardo Solar is a Corredanan \emph{sociedad anónima} (SA). Case P's Bélanou Power is a Kessaran \emph{société anonyme}, also SA. In the United States the usual form is a limited liability company: Sabine Pass Liquefaction, LLC, which borrowed for the first LNG export trains in Louisiana (\cref{ssec:2.3.3}), is one. In England it is a private limited company. The form matters less than three properties every form must have. The company's liability must be separate from its shareholders', its shares must be capable of being pledged to lenders, and it must be able to hold the licenses and sign the contracts the project needs. From here on the book calls the SPV the \term{project company}: the special purpose vehicle that owns the project, signs its contracts, and borrows its debt.

The owners of the project company are its sponsors. A \term{sponsor} is a company that develops a project and owns equity in its project company, carrying the first loss. Llano Pardo had two: Tallisford Energy Partners, a London fund manager, with 60\%, and Grupo Arismendi, a Corredanan family group, with 40\%. \Cref{ch:4} sorts sponsors into strategic, financial, and developer types and shows how each behaves; for now the definition needs only its last clause. The sponsors' equity is paid in before the lenders lose a dollar and is repaid only after the lenders have been paid what is due, so the sponsors are the first to absorb any shortfall.

Why form a new company at all, when Grupo Arismendi already had companies that could have signed the PPA? Each reason comes from a different chair at the table. The lenders want a borrower whose only assets and liabilities are the project's, so that the cash they forecast is the cash they will get, and so that no unrelated creditor stands beside them. The sponsors want their other businesses protected if the plant fails, and they want to be able to sell the project later by selling shares rather than by assigning dozens of contracts and licenses one by one. Tallisford, a fund with investors of its own, needed a holding it could value and eventually exit. ELNACOR and the Ministry of Energy and Mines wanted a counterparty whose sole business was delivering the plant they had tendered for, which made the PPA's performance obligations easier to enforce and the plant easier to keep running if a sponsor failed.

\subsection{What the ring-fence keeps in and keeps out}\label{ssec:2.1.2}

A new company is not enough on its own. Nothing in Corredanan company law stops Llano Pardo Solar SA from borrowing more money, guaranteeing a sponsor's bank loan, buying a second site, or paying its cash into an account Grupo Arismendi also uses. Each of those acts would let other creditors or other businesses reach the cash that the lenders sized their loan against. The finance documents therefore build a fence around the project company with promises.

\term{Ring-fencing} is the legal and contractual separation that keeps a project's assets and cash flows apart from its sponsors' other businesses, protecting each from the other's failure. It works in two directions, and practitioners who remember only one of them misjudge deals.

The first direction protects the sponsors from the project. If Llano Pardo had failed, Sterrenberg Bank and the Atlantic Basin Development Bank (ABDB) could have enforced their security over the plant, its contracts, its accounts, and the shares in the project company, and nothing else. Tallisford's other funds and Grupo Arismendi's agribusiness were beyond their reach. The sponsors' maximum loss was the equity they had put in.

The second direction protects the project's lenders from the sponsors. If Grupo Arismendi's agribusiness had collapsed in 2020, its creditors could have claimed its assets, including its 40\% shareholding in Llano Pardo Solar SA. Those shares were already pledged to the project lenders, so the agribusiness creditors would have taken them subject to the pledge, and the value they could extract was the shares' value after the project debt: in practice, the right to future distributions once the lenders had been paid. They could not seize the plant, the PPA receivables, or the cash in the project accounts, because none of those belonged to Grupo Arismendi. \Cref{exh:2.1} draws both directions on Llano Pardo's parties.

\begin{exhibit}[H]
\caption{The two directions of ring-fencing \illustrative}\label{exh:2.1}
\centering
\begin{tikzpicture}[node distance=10mm and 9mm]
\node[pfspv, text width=34mm] (pc) {Llano Pardo Solar SA\\(plant, PPA, accounts)};
\node[pfbox, above left=14mm and -4mm of pc, text width=30mm] (tal) {Tallisford fund\\(60\%)};
\node[pfbox, above right=14mm and -4mm of pc, text width=30mm] (ari) {Grupo Arismendi\\(40\%)};
\node[pfbox, above=of ari, text width=30mm] (aricr) {Grupo Arismendi's own creditors};
\node[pflend, below=14mm of pc, text width=34mm] (lend) {Project lenders\\(Sterrenberg Bank, ABDB)};
\node[pfbox, left=of lend, text width=24mm] (epc) {Montajes Cordillera};
\node[pfbox, right=of lend, text width=24mm] (off) {ELNACOR};
\draw[pfflow] (tal) -- node[pflabel] {Equity} (pc);
\draw[pfflow] (ari) -- node[pflabel] {Equity} (pc);
\draw[pfflow] (lend) -- node[pflabel] {Loan} (pc);
\draw[pfarrow] (pc) -- node[pflabel] {EPC contract} (epc);
\draw[pfarrow] (pc) -- node[pflabel] {PPA} (off);
\draw[pfarrow, dashed] (lend.west) to[out=150, in=200] node[pflabel, pos=0.55, text width=24mm, align=center] {No claim beyond the equity} (tal.west);
\draw[pfarrow, dashed] (aricr.east) to[out=-30, in=20] node[pflabel, pos=0.55, text width=22mm, align=center] {Shares only, subject to pledge} (pc.east);
\end{tikzpicture}
\exhibitsource{Chapter~1 deal specification; dashed arrows show the limits of each group's claim.}
\end{exhibit}

Four promises in the finance documents build the fence, and they usually sit together in the covenant section of the common terms agreement, the master finance document that all the senior lenders sign (\cref{ch:51} teaches the document set). The \term{single-purpose undertaking} is a covenant by the project company to carry on no business other than the project. A separateness undertaking, which the book uses without a formal definition, is the project company's promise to behave as a separate company: its own name, books, bank accounts, and decisions, and arm's-length terms with its owners. A limit on other financial debt and guarantees keeps new creditors from standing beside the lenders. A ban on mergers, acquisitions, and subsidiaries stops the project company from becoming something else. \Cref{cl:2.1} is an illustrative version drafted under English law for a 2024 common terms agreement for a wind farm in Chile's Biobío region, the deal that \cref{ex:2.1,ex:2.3,ex:2.4} follow.

\begin{clause}{Single-purpose and separateness undertaking, common terms agreement \illustrative}\label{cl:2.1}
\begin{enumerate}[label=(\alph*)]
\item \emph{Business.} The Borrower shall not carry on any business other than the development, ownership, construction, operation, maintenance and financing of the Project in accordance with the Project Documents, and activities reasonably incidental to it.
\item \emph{Financial Indebtedness and guarantees.} The Borrower shall not incur or permit to subsist any Financial Indebtedness, or give any guarantee or indemnity in respect of the obligations of any person, other than Permitted Financial Indebtedness.
\item \emph{Separateness.} The Borrower shall: (i) maintain its own books, records and bank accounts, separate from those of any Sponsor or Affiliate; (ii) not commingle its funds with those of any other person; (iii) hold itself out as a separate legal entity and conduct business in its own name; and (iv) enter into any transaction with a Sponsor or Affiliate only on arm's length terms, in writing, and, where the value of the transaction exceeds USD 250,000 in any Financial Year, with the prior consent of the Intercreditor Agent.
\item \emph{Corporate structure.} The Borrower shall not enter into any amalgamation, demerger, merger or corporate reconstruction, acquire any company, business or undertaking, or establish or acquire any Subsidiary.
\end{enumerate}
\end{clause}
\begin{enumerate}
\item Paragraph (a). The words ``reasonably incidental'' carry the negotiation. Sponsors want them wide enough to cover future grid services or a co-located battery without asking; lenders want them narrow, because each new activity brings new risks the model did not test. The usual landing is a narrow clause plus a consent process.
\item Paragraph (b). ``Permitted Financial Indebtedness'' is a defined list (typically the senior facilities, hedging under the hedging policy, subordinated shareholder loans, and a small working-capital basket). The guarantee ban protects the second direction of the ring-fence: without it, a sponsor could ask the project company to guarantee the sponsor's own corporate loan.
\item Paragraph (c). Separateness is what persuades a court in a sponsor's insolvency that the project company is a different person. A project company that shared bank accounts with its parent, paid the parent's bills, and let the parent sign for it invites a creditor's argument that the two should be treated as one. The consent threshold in sub-paragraph (iv) aims at the affiliate EPC or O\&M contract on soft terms, the most common way value leaks out of a fenced project.
\item Paragraph (d). A project company that merged with another company would import that company's creditors. Lenders sometimes allow a subsidiary for a specific purpose, such as holding a transmission line that must later be transferred to a utility, and list it as an exception.
\end{enumerate}

Whether the fence holds in the end is a question for the insolvency law of the sponsor's country, not for the clause. Common-law courts can, in rare cases, consolidate a parent and a subsidiary that were run as one business; civil-law systems reach similar results through other doctrines. The clause cannot prevent that, but a project company that has actually behaved separately gives a court little reason to ignore its separateness. \Cref{ch:52} shows how the security package and the accounts structure reinforce the fence, and \cref{ch:64} what happens to it in a sponsor's insolvency.

\subsection{SunEdison, TerraForm, and a ring-fence under test}\label{ssec:2.1.3}

SunEdison's collapse tested both directions of the ring-fence at once, and the result was mixed in a way that teaches more than a clean success would have.

The structure had three layers. At the bottom sat individual wind and solar projects, many with their own non-recourse project debt. In the middle sat TerraForm Power, a listed company that SunEdison had floated on NASDAQ in July 2014 to buy operating projects from it. TerraForm Power borrowed on its own account: in 2015 its subsidiary Terra Operating LLC issued USD~800~million of 5.875\% senior notes due 2023 and USD~300~million of 6.125\% senior notes due 2025 to pay for acquisitions. At the top sat SunEdison, which developed new projects, dropped finished ones into TerraForm Power, and kept control through high-vote Class~B shares and incentive distribution rights, a claim to a rising share of TerraForm Power's distributions (TerraForm Power 2016). \Cref{ch:32} teaches this listed-vehicle structure, known as a yieldco, in full; here only its fence matters.

The fence around the projects and around TerraForm Power's cash held. Neither TerraForm company filed for bankruptcy, neither was consolidated into SunEdison's case, and the project-level lenders kept their claims on their own projects. When Brookfield agreed in March 2017 to become TerraForm Power's sponsor, it bought into a business that had stayed whole.

The fence around governance did not hold, because SunEdison's control rights crossed it by design. By January 2016 SunEdison held 84.1\% of TerraForm Power's votes against 34.5\% of its economic interest. On November 20, 2015, it used that control to remove TerraForm Power's chief executive and chief financial officer and to replace members of the committee meant to police conflicts with SunEdison; three directors resigned in protest. SunEdison had also pledged its Class~B shares and its incentive distribution rights to its own lenders, so in bankruptcy SunEdison's creditors held the levers of control over TerraForm Power. TerraForm Power's 2015 accounts were filed on December 5, 2016, with a going-concern note citing the risk of substantive consolidation, and its late filings put the reporting covenants in its note indentures and its NASDAQ listing at risk. SunEdison's creditors' committee pursued claims to unwind payments and asset transfers between SunEdison and TerraForm Power, released only in the 2017 settlement with Brookfield (TerraForm Power 2016; TerraForm Power 2018).

The mechanism to carry away is that a ring-fence is a set of contracts and corporate separations, and it protects exactly what those contracts and separations cover. Debt and assets were fenced. Votes, management, services, and the pipeline of new projects were deliberately connected to the sponsor, so the sponsor's failure flowed through them. When you review a structure, ask three questions beyond the debt: who appoints the board, who provides the services the company cannot run without, and which agreements fall away if the sponsor fails. \Cref{ssec:14.18.1} returns to SunEdison when it teaches sponsor credit risk.

